% Ad Atticum 11.3 (ad-atticum-11-03) --- English translation
% Translated by Claude (Anthropic) from the Latin (Perseus).
% Style guide: STYLE.md.

\ciceroLetterOpener{CICERO ATTICO salutem}

\ciceroSection{1} What is going on here you will be able to learn from the man who brings this letter. I have kept him longer than I should because we have been waiting day by day for some news; and even now there has been no reason to send him on except the matter on which you wished a reply, namely what I want done with regard to the Kalends of Quintilis. Both choices are heavy: at so heavy a time the risk of so great a sum, and with the outcome of affairs in the balance, that breaking off you write of. So, here as elsewhere, this above all I commit to your care and goodwill, and to her judgement and her wishes; for whom in her wretchedness I should have provided better, if, while there was still time, I had taken counsel with you face to face rather than by letter about my survival and my fortunes.

\ciceroSection{2} As to your denial that any peculiar trouble of my own hangs over me in the common misfortunes, although there is some consolation in your saying so, even so there are many particular troubles which you surely see are most grievous and which I could have most easily avoided. They will, however, be less if, as has been the case so far, they are lightened by your management and your care.

\ciceroSection{3} The money is with Egnatius. Let it remain with me where it is. For what is going on here does not look as if it can last long, so that I may yet be able to know what is most needed --- though I am in want of everything, since the man I am with is also in straits: I lent him a large sum, supposing that, once affairs were settled, the loan would even be a credit to me. As before, please draft yourself whatever letters you judge ought to be written by me to anyone. Greet your people for me. Take care to keep well. Above all, with regard to what you write, in all things look after and provide that nothing be wanting to the one for whose sake you know me to be most miserable. On the Ides of June, from the camp.
